\begin{table}[htbp]
\centering
\begin{threeparttable}
\caption{Preg\~ao drop-out costs $F_c^k$ by stratum (point ID)}
\label{tab:v3_pregao_fc_summary}
\small
\begin{tabular}{llllrrrrr}
\toprule
Class & Type & Period & N losers & Mean & p50 & p75 & p90 & \% $>$ ref \\
\midrule
non-pharma & non-SME & Post & 26,340 & 0.9977 & 0.8636 & 1.2500 & 1.8257 & 38.52 \\
non-pharma & non-SME & Pre & 43,746 & 0.9603 & 0.8202 & 1.2174 & 1.8010 & 36.07 \\
non-pharma & SME & Post & 33,543 & 1.1007 & 0.9606 & 1.3978 & 2.0299 & 46.42 \\
non-pharma & SME & Pre & 16,701 & 1.1069 & 0.9537 & 1.4355 & 2.0526 & 46.49 \\
pharma & non-SME & Post & 26,928 & 0.8591 & 0.7677 & 1.0615 & 1.4981 & 29.41 \\
pharma & non-SME & Pre & 43,668 & 0.7915 & 0.7038 & 1.0000 & 1.4125 & 24.47 \\
pharma & SME & Post & 17,491 & 1.0606 & 0.9502 & 1.3883 & 2.0000 & 45.90 \\
pharma & SME & Pre & 8,063 & 1.0471 & 0.8590 & 1.4575 & 2.2321 & 41.16 \\
\bottomrule
\end{tabular}
\begin{tablenotes}\footnotesize
\item Normalized by item reference price; values near 1 mean the
firm bid close to the ceiling. Under English-reverse IPV, each loser
drops out at its cost, so the ECDF of losers' drop-outs is a point
estimate of $F_c^k$.
\end{tablenotes}
\end{threeparttable}
\end{table}
